% !TEX root = projeto.tex
% Revisado Renato 31 de agosto - versão 1.0
\setlength{\absparsep}{18pt}
\begin{resumo}[Resumo]
Neste projeto de TCC I, apresento a fundamentação e o planejamento da evolução do Estoque Fácil, protótipo web que venho desenvolvendo para gestão e controle de estoques. Defini o tema a partir de observações exploratórias realizadas durante minha atuação no suporte a um ERP em Sinop, sem tratá-las como levantamento representativo. Parto da base de código já desenvolvida e de demandas documentadas do domínio, com foco no controle de materiais em pequenas e médias organizações. Analiso como requisitos, tecnologias, arquitetura, infraestrutura e critérios de qualidade podem orientar a evolução do protótipo. Delimito o problema, os objetivos, a metodologia, o objeto de estudo e o protocolo de avaliação. Neste TCC I, não apresento o Estoque Fácil como produto concluído nem antecipo resultados de testes ainda não executados. No TCC II, pretendo implementar as complementações priorizadas e avaliar uma versão identificada do sistema.

\textbf{Palavras-chave}: gestão de estoques; sistema web; engenharia de software; arquitetura de software; qualidade de software.
\end{resumo}

\begin{resumo}[Abstract]
In this first stage of my undergraduate thesis (TCC I), I present the theoretical foundation and plan for evolving Estoque Fácil, a web prototype I have been developing for inventory management and control. I selected the topic from exploratory observations made while providing support for an ERP system in Sinop, Mato Grosso, without treating them as a representative survey of local businesses. I use the existing code base and documented domain needs to examine how requirements, technologies, architecture, infrastructure, and software quality criteria can guide the prototype's evolution. I delimit the research problem, objectives, methodology, object of study, and evaluation protocol. At this stage, I do not present Estoque Fácil as a completed product or claim results from tests that have not yet been performed. In TCC II, I intend to implement the prioritized additions and evaluate an identified version of the system.

\textbf{Keywords}: inventory management; web system; software engineering; software architecture; software quality.
\end{resumo}
